\section{Conclusion}
\label{sec:conclusion}

This study provides a comprehensive global characterization of the Martian crustal magnetic field and evaluates its biological and operational implications across active and candidate landing sites. Our analysis establishes that candidate human landing sites located in the northern plains (Arcadia Planitia, Deuteronilus Mensae) and rover sites (Jezero Crater, Oxia Planum) reside within extreme near-null hypomagnetic environments ($|\mathbf{B}_{\text{surf}}| < \SI{20}{\nano\tesla}$), representing a $>2,500$-fold reduction relative to Earth's geomagnetic baseline. In contrast, the southern highlands preserve intense crustal anomalies producing localized mini-magnetospheres exceeding \SI{1500}{\nano\tesla} at ground level. Synthesizing these physical constraints with terrestrial magnetobiology literature indicates that unmitigated hypomagnetic exposure poses distinct developmental risks to crop plants, alters closed-loop microbial dynamics, and may synergistically accelerate musculoskeletal bone loss in astronauts. We emphasize the necessity of deploying surface magnetometers on upcoming landers and recommend incorporating artificial magnetic field compensation into long-duration Martian agricultural and habitat growth modules.
